\section*{NeurIPS Paper Checklist}

\begin{enumerate}

\item {\bf Claims}
    \item[] Question: Do the main claims made in the abstract and introduction accurately reflect the paper's contributions and scope?
    \item[] Answer: \answerYes{}
    \item[] Justification: The abstract and introduction clearly state the three main contributions (state representation design, truncation-aware GAE, inference-time ensemble) and report the corresponding quantitative results (mean reward, programs beating \texttt{-Oz}). All claims are directly supported by experiments in Section~4, including a cross-architecture evaluation (Section~4.5) where the proposed \texttt{pass\_sequence} feature is integrated into a recurrent GRU baseline, demonstrating its architectural robustness and independently replicating the ablation finding across two systems developed in parallel.
    
\item {\bf Limitations}
    \item[] Question: Does the paper discuss the limitations of the work performed by the authors?
    \item[] Answer: \answerYes{}
    \item[] Justification: Section~5 (Conclusion) includes a dedicated Limitations paragraph discussing the absence of GA/SA/MCTS baselines under controlled compute budget, the structural failure of the 30-pass action set on \texttt{linux-v0} kernel code, and the saturation of the current PPO design space.

\item {\bf Theory assumptions and proofs}
    \item[] Question: For each theoretical result, does the paper provide the full set of assumptions and a complete (and correct) proof?
    \item[] Answer: \answerNA{}
    \item[] Justification: This paper presents no theoretical results; all contributions are empirical.

\item {\bf Experimental result reproducibility}
    \item[] Question: Does the paper fully disclose all the information needed to reproduce the main experimental results of the paper to the extent that it affects the main claims and/or conclusions of the paper (regardless of whether the code and data are provided or not)?
    \item[] Answer: \answerYes{}
    \item[] Justification: Section~4.1 specifies all PPO hyperparameters ($\varepsilon{=}0.2$, $\gamma{=}0.99$, $\lambda{=}0.95$, n\_steps, entropy coefficient, network architecture), the action set, reward space, episode length, and number of training episodes for both Phase~1 (MLP) and Phase~2 (GRU) architectures. The environment (CompilerGym \texttt{llvm-v0}, \texttt{cbench-v1}) is publicly available under the Apache 2.0 license.

\item {\bf Open access to data and code}
    \item[] Question: Does the paper provide open access to the data and code, with sufficient instructions to faithfully reproduce the main experimental results, as described in supplemental material?
    \item[] Answer: \answerNo{}
    \item[] Justification: Code is not released with this submission due to double-blind anonymity requirements. The benchmark data (\texttt{cbench-v1}) is publicly available through CompilerGym~\cite{CompilerGym}. All hyperparameters, network architectures, and implementation details required for reproduction are fully specified in Section~4.1 and the PPO configuration table. Code will be released upon acceptance.

\item {\bf Experimental setting/details}
    \item[] Question: Does the paper specify all the training and test details (e.g., data splits, hyperparameters, how they were chosen, type of optimizer) necessary to understand the results?
    \item[] Answer: \answerYes{}
    \item[] Justification: Section~4.1 provides full training details for both architectures: Phase~1 uses a feed-forward 256$\times$256 MLP over the full 124-pass action space (500 episodes per benchmark, NVIDIA RTX 5070 Ti), and Phase~2 uses a recurrent GRU over the curated 30-pass action space (300k steps, 8 parallel workers, Apple Mac mini M4 Pro). Hyperparameters were chosen via the systematic ablation studies in Sections~4.2--4.4, covering sequence window $k$, network capacity, n\_steps, and entropy coefficient.

\item {\bf Experiment statistical significance}
    \item[] Question: Does the paper report error bars suitably and correctly defined or other appropriate information about the statistical significance of the experiments?
    \item[] Answer: \answerNo{}
    \item[] Justification: Error bars over multiple training seeds are omitted due to the high computational cost of the compiler environment. However, to establish statistical significance, Section~4.6 introduces a compute-budget-matched baseline (Random best-of-56) that gives the random policy $3.5\times$ more rollouts than our PPO (56 vs.\ 16). Our learned policy outperforms this stronger baseline on all 23 benchmarks (avg 1.060 vs.\ 0.992), providing strong evidence that gains stem from policy knowledge rather than sampling variance. The ablation study further evaluates each configuration across all 23 \texttt{cbench-v1} benchmarks, and the count of programs beating \texttt{-Oz} provides a distribution-free summary statistic.

\item {\bf Experiments compute resources}
    \item[] Question: For each experiment, does the paper provide sufficient information on the computer resources (type of compute workers, memory, time of execution) needed to reproduce the experiments?
    \item[] Answer: \answerYes{}
    \item[] Justification: Phase~1 ablation experiments (Sections~4.2--4.4) and the compute-budget comparison (Section~4.6) were conducted on a machine with an \textbf{NVIDIA RTX 5070 Ti} GPU (500 episodes per benchmark). Phase~2 ensemble training and cross-suite evaluation (Sections~4.7--4.8) were conducted on an \textbf{Apple Mac mini M4 Pro} (24\,GB unified memory) using 8 parallel CPU workers, where each 300k-step training run completes in approximately 47 minutes. Cross-architecture validation (Section~4.5) uses the Phase~2 hardware for GRU-Base、GRU-Seq、GRU-Seq-Full.

\item {\bf Code of ethics}
    \item[] Question: Does the research conducted in the paper conform, in every respect, with the NeurIPS Code of Ethics?
    \item[] Answer: \answerYes{}
    \item[] Justification: This work studies compiler optimization, a purely technical domain with no human subjects, sensitive data, or dual-use concerns. The research conforms with the NeurIPS Code of Ethics.

\item {\bf Broader impacts}
    \item[] Question: Does the paper discuss both potential positive societal impacts and negative societal impacts of the work performed?
    \item[] Answer: \answerNA{}
    \item[] Justification: This work improves compiler optimization efficiency, which has broadly positive effects (reduced energy consumption, faster software). There are no foreseeable negative societal impacts from automated compiler pass ordering.

\item {\bf Safeguards}
    \item[] Question: Does the paper describe safeguards that have been put in place for responsible release of data or models that have a high risk for misuse (e.g., pre-trained language models, image generators, or scraped datasets)?
    \item[] Answer: \answerNA{}
    \item[] Justification: This paper poses no misuse risks. The trained models are compiler optimization policies operating solely within the LLVM compiler infrastructure, with no dual-use potential.

\item {\bf Licenses for existing assets}
    \item[] Question: Are the creators or original owners of assets (e.g., code, data, models), used in the paper, properly credited and are the license and terms of use explicitly mentioned and properly respected?
    \item[] Answer: \answerYes{}
    \item[] Justification: CompilerGym~\cite{CompilerGym} and the \texttt{cbench-v1} benchmark suite are cited in the paper and used in accordance with the Apache 2.0 license. All other third-party dependencies (PyTorch, Stable-Baselines3) are similarly open-source and used within their respective license terms.

\item {\bf New assets}
    \item[] Question: Are new assets introduced in the paper well documented and is the documentation provided alongside the assets?
    \item[] Answer: \answerNA{}
    \item[] Justification: This paper does not release new datasets or pretrained model weights. The primary contributions are algorithmic: the \texttt{pass\_sequence} state representation, the truncation-aware GAE implementation, and the inference-time ensemble strategy. Implementation details sufficient for re-implementation are provided in Sections~3 and~4.

\item {\bf Crowdsourcing and research with human subjects}
    \item[] Question: For crowdsourcing experiments and research with human subjects, does the paper include the full text of instructions given to participants and screenshots, if applicable, as well as details about compensation (if any)?
    \item[] Answer: \answerNA{}
    \item[] Justification: This paper does not involve crowdsourcing or human subjects.

\item {\bf Institutional review board (IRB) approvals or equivalent for research with human subjects}
    \item[] Question: Does the paper describe potential risks incurred by study participants, whether such risks were disclosed to the subjects, and whether Institutional Review Board (IRB) approvals (or an equivalent approval/review based on the requirements of your country or institution) were obtained?
    \item[] Answer: \answerNA{}
    \item[] Justification: This paper does not involve human subjects research.

\item {\bf Declaration of LLM usage}
    \item[] Question: Does the paper describe the usage of LLMs if it is an important, original, or non-standard component of the core methods in this research?
    \item[] Answer: \answerNA{}
    \item[] Justification: LLMs are not used in the core methodology of this work. The paper explicitly positions itself as a traditional DRL alternative to LLM-based compiler optimization approaches (Section~2.2), and all policy networks are conventional MLP and GRU architectures trained with PPO.

\end{enumerate}